% LaTeX companion of 08-backward-error-analysis.typ. Both formats are editable.
\chapter{backward error analysis of NLA algorithms}

Source attribution in the selected TeX chapter title: \texttt{doi:10.1137/1.9780898719574.ch3}.

我们在上一个 Ch 中介绍了 conditioning 和 stability. 现在我们用它们对经典的 NLA algorithms 进行 backward error analysis.

\section{Stability of Householder Triangularization}

我们 set \(R,Q\) to be random upper triangular 和 orthogonal matrices (by orthogonizing 一个 random matrix)，并 set \(A := QR\).

\begin{verbatim}
R = triu (randn(50));
[Q,X] = qr(randn(50));
A = Q*R
\end{verbatim}

然后我们再对 \(A\) 进行 QR 分解, via Household (Matlab 自带使用 Household), 看看 relative error:

\begin{verbatim}
[Q2,R2] = qr(A);
norm (Q2 - Q);
    ans = 0.00889
norm (R2-R) / norm(R);
    ans = 0.00071
\end{verbatim}

我们发现 \(Q,R\) 的 relative error 其实很大.

但是，当我们用这个 \(QR\) 计算 \(A\) 时:

\begin{verbatim}
norm (A - Q2*R2) / norm(A);
    ans = 1.432e-15
\end{verbatim}

我们发现一个惊人的事实: 这个 QR 分解的 error

\(\text{Q3} = Q + 1\ \text{e} - 4 \ast \text{randn}(50)\)

\(\text{R3} = R + 1\ \text{e} - 4 \ast \text{randn}(50)\)

\(\frac{\text{norm}\left( {A - \ \text{Q3}\  \ast \ \text{R3}} \right)}{\text{norm}(A)}\)

\(\text{ans} = 0.00088\)

\section{Stability of Back Substitution}

\section{Conditioning of Least Squares}

\section{Stability of Least Squares}
